% ============================================================================
%  Manuscrito para Astronomy & Astrophysics
%  Clase aa.cls v9.4 (EDP Sciences, 2026) — ver class/
%
%  Compilar desde este directorio:
%      latexmk -pdf main.tex
%  o, si no hay TeX local, subir build/submission/ a Overleaf (ver README.md).
%
%  ANTES DE ENVIAR:  uv run python tools/flatten.py
% ============================================================================

% -- Opciones de clase -------------------------------------------------------
%   (ninguna)   layout estándar A&A, dos columnas
%   letter      Letter to the Editor (la mención se añade sola)
%   longauth    listas muy largas de autores/afiliaciones (van tras las refs.)
%   onecolumn   artículos con fórmulas grandes; lo confirma el editor
%   bibyear     si NO usas referencias estructuradas autor-año de natbib
%   referee     versión para árbitro
% ----------------------------------------------------------------------------
\documentclass{aa}

% A&A usa fuentes PostScript TX Times, no Computer Modern.
\usepackage[varg]{txfonts}

\usepackage[utf8]{inputenc}
\usepackage[T1]{fontenc}

% -- Matemáticas -------------------------------------------------------------
\usepackage{amsmath}
\usepackage{amssymb}

% -- Figuras y floats --------------------------------------------------------
\usepackage{graphicx}
\usepackage{subcaption}   % necesario para \ContinuedFloat
\usepackage{lscape}       % tablas/figuras rotadas a página completa
\usepackage{placeins}     % \FloatBarrier

% -- Tablas ------------------------------------------------------------------
\usepackage{longtable}
\usepackage{csvsimple}
\usepackage{xparse}
\usepackage{etoolbox}

% -- Bibliografía ------------------------------------------------------------
\usepackage{natbib}
\bibpunct{(}{)}{;}{a}{}{,}   % puntuación exigida por el estilo A&A

% -- Comandos propios de autoría --------------------------------------------
%  OJO: A&A pide no usar definiciones propias en el archivo enviado.
%  tools/flatten.py las expande antes del envío. Ver aapaper.sty.
\usepackage{aapaper}

\graphicspath{{figures/}}

\begin{document}

% ============================================================================
%  CABECERA: título, autores, afiliaciones, abstract, keywords
% ============================================================================
% ============================================================================
%  Cabecera del manuscrito: título, autores, afiliaciones, fecha
%  El abstract y las keywords viven en sections/00-abstract.tex
%
%  Reglas A&A (manuscript header):
%    - Separar cada autor con \and.
%    - \corrauth{correo} marca al autor de correspondencia; aa.cls v9.4 lo
%      convierte solo en una nota al pie. Un único \corrauth por artículo.
%    - \email{correo} NO se imprime: lo extrae el sistema editorial de EDPS
%      para que cada coautor autentique su ORCID.
%    - NO poner ORCIDs junto a los nombres. Se eliminan en producción.
%    - El título puede partirse con \\ y llevar \thanks (separados con \fnmsep).
% ============================================================================

\title{Título del artículo\thanks{Nota al pie del título, si hace falta:
       agencia de financiación, número de proyecto, etc.}}

% Quitar por completo si no hay subtítulo.
\subtitle{Subtítulo}

\author{
  Luis Daniel Díaz\inst{1}\corrauth{correo@udea.edu.co}
  \and Nombre Coautor\inst{1}\email{coautor@udea.edu.co}
  \and Nombre Coautor\inst{2}\email{coautor@otra.edu}
}

\institute{
  Instituto de Física, Universidad de Antioquia, Calle 70 No. 52-21,
  Medellín, Colombia
  \and
  Segunda afiliación, Ciudad, País
}

\date{Received DD Month YYYY / Accepted DD Month YYYY}

% ============================================================================
%  ABSTRACT Y KEYWORDS
%
%  A&A usa abstract ESTRUCTURADO: \abstract toma CINCO argumentos entre llaves,
%  y los cinco son obligatorios como tokens, aunque el 1.o y el 5.o puedan ir
%  vacíos ({}).
%
%      1. Context      opcional  — trasfondo del problema
%      2. Aims         OBLIGATORIO — qué se propone el artículo
%      3. Methods      OBLIGATORIO — cómo se investigó
%      4. Results      OBLIGATORIO — qué se obtuvo
%      5. Conclusions  opcional  — qué se concluye en general
%
%  Los encabezados ("Context.", "Aims.", ...) los pone la clase; no escribirlos.
%
%  Si prefieres el abstract tradicional de un solo párrafo, usa la forma de un
%  único argumento: \abstract{...} — pero la estructurada es la preferida.
%
%  NO dejar líneas en blanco ni comentarios ENTRE los cinco grupos: una línea
%  vacía abre un párrafo nuevo y rompe la lectura de los argumentos. Por eso
%  van pegados abajo, y toda la explicación vive en este bloque.
%
%  Toda abreviatura de instrumento, método u observatorio se expande la primera
%  vez aquí y otra vez en la introducción:
%      ...very long baseline interferometry (VLBI)...
% ============================================================================

\abstract
{}
{Objetivo del trabajo.}
{Métodos empleados.}
{Resultados obtenidos.}
{}

% ----------------------------------------------------------------------------
%  KEYWORDS
%  Separadas por " -- " (en-dash rodeado de espacios).
%  Tomarlas TEXTUALMENTE de la lista oficial de A&A; no inventar.
%  Lista: https://www.aanda.org/for-authors/latex-issues/information-files
%
%  Candidatas para este proyecto (verificar en la lista antes de fijarlas):
%    black hole physics
%    gravitation
%    relativistic processes
%    celestial mechanics
%    methods: numerical
%    Galaxy: center
% ----------------------------------------------------------------------------
\keywords{black hole physics --
          gravitation --
          methods: numerical}


\maketitle

% Descomentar solo para publicar en arXiv (quita la numeración de líneas):
% \nolinenumbers


% ============================================================================
%  CUERPO
%  Orden fijo por prefijo numérico. Para reordenar, renombra los archivos.
% ============================================================================
% ============================================================================
%  1. INTRODUCTION
%
%  A&A exige que el manuscrito empiece por "1. Introduction". La numeración la
%  lleva LaTeX: nunca escribir el número a mano.
%
%  Jerarquía disponible:
%      \section{Título}        -> 1.
%      \subsection{Título}     -> 1.1
%      \subsubsection{Título}  -> 1.1.1
%      \paragraph{Título}
%  Toda sección necesita un título corto y descriptivo.
%
%  Cross-referencing: \label en TODA sección, figura, tabla y ecuación
%  referenciada; \ref para citarlas. A&A lo convierte en hipervínculos HTML.
%  Atajos de aapaper.sty: \secref{} \figref{} \tabref{} \eqnref{} \appref{}
%
%  Objetos astronómicos: envolverlos SIEMPRE en \object{}, también dentro de
%  tablas pequeñas. Genera el archivo .obj que enlaza con SIMBAD/ALADIN.
%      \object{Sgr A*}   \object{S2}
%
%  Abreviaturas: expandir la primera vez que aparecen en el texto principal.
% ============================================================================

\section{Introduction}
\label{sec:introduction}

Texto de la introducción \citep{baker66}. Esa cita es un placeholder: un
\texttt{thebibliography} sin ninguna entrada es un error de LaTeX, así que el
esqueleto trae una para compilar. Bórrala junto con la entrada de
\texttt{refs.bib} en cuanto tengas tus propias referencias.

% ----------------------------------------------------------------------------
%  Ejemplos de uso — descomentar y adaptar
% ----------------------------------------------------------------------------
%
%  Citas (natbib con la puntuación A&A ya configurada en main.tex):
%      \citet{gillessen17}          -> Gillessen et al. (2017)
%      \citep{gillessen17}          -> (Gillessen et al. 2017)
%      \citep[see][chap.~2]{do19}   -> (see Do et al. 2019, chap. 2)
%      \citep{a18,b19}              -> (A et al. 2018; B et al. 2019)
%
%  Referencias cruzadas:
%      as shown in \secref{sec:methods}
%      \figref{fig:orbit} displays the projected orbit
%      \tabref{tab:elements} lists the orbital elements
%      from \eqnref{eq:geodesic}
%
%  Objetos:
%      The star \object{S2} orbits \object{Sgr A*} with a period of 16 yr.
%
%  Tipografía A&A (ver el checklist completo en la skill aa-paper):
%      20\,000 km              miles con espacio fino
%      1950--1985              rango con en-dash
%      $-30$~K                 menos matemático + espacio fijo
%      $v = 7650$~km~s$^{-1}$  unidad con índice negativo, nunca km/s
%      $T_\mathrm{eff}$        subíndice-etiqueta en redonda
%      \ion{H}{ii}             ion, romano en minúsculas
% ----------------------------------------------------------------------------

% ============================================================================
%  2. METHODS
%
%  Reglas A&A sobre ecuaciones:
%    - Toda ecuación referenciada con \ref necesita su \label. Solo ese
%      mecanismo; nada de numerar a mano.
%    - Puntuar las ecuaciones como si fueran texto: coma o punto al final.
%    - Delimitadores con \left( ... \right) para que escalen con la fórmula.
%    - En modo matemático, todo va en itálica por defecto. Lo que NO es una
%      variable va en redonda:
%          \mathrm{d}   diferencial
%          \mathrm{e}   base del logaritmo natural
%          \mathrm{i}   unidad imaginaria
%          $T_\mathrm{eff}$, $m_\mathrm{e}$   subíndices que son etiquetas
%          \mbox{Hz}, \mathrm{km\,s^{-1}}     unidades
%      Pero SÍ en itálica si el subíndice es a su vez una variable.
%    - Vectores: \vec{A}.  Tensores: \tens{ABC}.  Solo en modo matemático.
%    - La numeración de ecuaciones en A&A es correlativa en todo el artículo,
%      no por sección.
% ============================================================================

\section{Methods}
\label{sec:methods}

Texto de la sección.

% ----------------------------------------------------------------------------
%  Ejemplo de ecuación numerada y referenciable
% ----------------------------------------------------------------------------
% \begin{equation}
%   \label{eq:geodesic}
%   \frac{\mathrm{d}^2 x^\mu}{\mathrm{d}\tau^2}
%     + \Gamma^\mu_{\ \alpha\beta}
%       \frac{\mathrm{d}x^\alpha}{\mathrm{d}\tau}
%       \frac{\mathrm{d}x^\beta}{\mathrm{d}\tau} = 0 \,,
% \end{equation}
%
% donde $\tau$ es el tiempo propio.

\subsection{Subsección}
\label{sec:methods-sub}

Texto.

% ============================================================================
%  3. RESULTS
%
%  Esta es la sección donde normalmente viven las figuras y las tablas.
%  Referencia rápida de los comandos de aapaper.sty; el detalle está en
%  la skill global aa-paper (references/figures.md y references/tables.md).
%
%  ------------------------------------------------------------------ FIGURAS
%  A&A admite tres anchos y solo tres:
%
%    \aafig[label]{archivo}{caption}[colocación]        88 mm, una columna
%    \aafigside[label]{archivo}{caption}[colocación]   120 mm, caption al lado
%    \aafigwide[label]{archivo}{caption}[colocación]   180 mm, dos columnas
%    \aafigmulti[label]{f1,f2,f3}{caption}[colocación] varios gráficos
%
%  El archivo se busca en figures/ (ver \graphicspath en main.tex).
%  Los floats a dos columnas hay que declararlos TEMPRANO en el fuente:
%  LaTeX solo puede empujarlos hacia adelante, nunca hacia atrás.
%
%  Requisitos de las imágenes:
%    - eps, pdf, jpg, png o tiff. Bitmaps a 250-300 dpi mínimo, no 92.
%    - Recortar el espacio vacío alrededor.
%    - Sin grids ni fondos de color salvo que aporten información.
%    - Nada de rojo+verde juntos (accesibilidad para daltonismo).
%    - Sin líneas finas; en minúsculas el texto dentro de la figura.
%    - Los símbolos se explican en el caption, no dentro de la figura.
%
%  ------------------------------------------------------------------- TABLAS
%    \aatable[label]{archivo.csv}{caption}[nota]        una columna
%    \aatablewide[label]{archivo.csv}{caption}[nota]    dos columnas
%    \aatablelong[n]{archivo.csv}{caption}{label}       multipágina
%    \aatablebib{(1)~\citet{...}; (2) \citet{...}.}     refs. bajo la tabla
%
%  El CSV va en tables/. Primera fila = encabezado. Las celdas admiten LaTeX
%  crudo ($M_\odot$, \tablefootmark{a}). Si un valor lleva coma, comillas
%  dobles alrededor.
%
%  Los comandos ya aplican las reglas A&A: caption arriba, \hline\hline /
%  \hline / \hline (NO booktabs), columnas a la izquierda, sin líneas
%  verticales. La tabla completa no debe exceder 23.5 cm de alto.
%
%  Notas al pie de tabla:
%    - general:     el 4.o argumento de \aatable pasa a \tablefoot{}
%    - referenciada: \tablefootmark{a} en la celda del CSV, y en la nota
%                    \tablefoottext{a}{texto de la nota}
% ============================================================================

\section{Results}
\label{sec:results}

Texto de la sección. Ejemplo vivo de los comandos: la \figref{fig:demo} y la
\tabref{tab:demo} de abajo compilan y se expanden con \texttt{tools/flatten.py}.
Bórralas cuando empieces a escribir de verdad.

\aafig[fig:demo]{demo}{Figura de demostración.}

\aatable[tab:demo]{tables/demo.csv}
  {Tabla de demostración generada desde \texttt{tables/demo.csv}.}
  [Nota general de la tabla.
   \tablefoottext{a}{Nota referenciada desde una celda del CSV.}]

% ----------------------------------------------------------------------------
%  Ejemplos — descomentar y adaptar
% ----------------------------------------------------------------------------
%
% \aafig[fig:orbit]{orbit.pdf}{Projected orbit of \object{S2} around
%   \object{Sgr A*}. Filled circles mark the astrometric epochs.}
%
% \aafigwide[fig:residuals]{residuals.pdf}{Residuals of the fit in right
%   ascension (\emph{left}) and declination (\emph{right}).}
%
% \aafigside[fig:posterior]{posterior.pdf}{Marginalised posterior distribution
%   of the black hole mass.}
%
% \aatable[tab:elements]{tables/elements.csv}
%   {Best-fitting orbital elements.}
%   [Uncertainties are the 68\% credible intervals.
%    \tablefoottext{a}{Held fixed during the fit.}]
%
% \aatablewide[tab:comparison]{tables/comparison.csv}
%   {Comparison with previous determinations.}
%
% \aatablelong[1]{tables/catalogue.csv}{Full astrometric catalogue.}{tab:catalogue}
% ----------------------------------------------------------------------------

% ============================================================================
%  4. DISCUSSION
%
%  Nota de estilo de A&A: la revista publica una lista de correcciones
%  idiomáticas frecuentes. Las más habituales en manuscritos de dinámica
%  relativista:
%
%    "depend on", no "depend of"          "independent of", no "independent on"
%    "associated with", no "associated to" "comparable to", no "comparable with"
%    "typical of", no "typical for"        "adjacent to", no "adjacent with"
%    "on the order of" (aprox.), no "of order"
%    "in the past 5 years", no "in the last 5 years"
%    "such as", no "like for example"      "we describe X in detail", no "in details"
%    "in the case of", no "in case of"     "on a timescale of", no "in a timescale"
%    "originate in", no "originate from"   "to stem from", mejor que "to result from"
%    "search for X", no "search of X"      "least-square technique" (con guion)
%    "Galactic bulge", "Galaxy" en mayúscula si son la nuestra
%    Evitar "quite / rather / somewhat": el adjetivo solo suele bastar.
%    "but" y "however" dicen lo mismo: usar uno de los dos, no ambos.
%
%  Alto/bajo vs. grande/pequeño vs. fuerte/débil:
%    high/low     -> value, rate, redshift, temperature, luminosity, velocity,
%                    energy, density, eccentricity, inclination, flux
%    large/small  -> scale, amplitude, mass, momentum, uncertainties
%    short/long   -> time, length, timescale
%    strong/weak  -> correlation, gradient, dependence, constraint, instability
% ============================================================================

\section{Discussion}
\label{sec:discussion}

Texto de la sección.

% ============================================================================
%  5. CONCLUSIONS
% ============================================================================

\section{Conclusions}
\label{sec:conclusions}

Texto de la sección.

% ============================================================================
%  AGRADECIMIENTOS
%
%  Entorno propio de aa.cls. Va justo ANTES de la bibliografía, sin \section:
%  la clase le da su propio formato. No lleva número.
%
%  Aquí van también las comunicaciones privadas, que nunca entran en la lista
%  de referencias:  "(Nombre YYYY, priv. comm.)"
% ============================================================================

\begin{acknowledgements}
  Texto de los agradecimientos: financiación, número de proyecto,
  infraestructura de cómputo, personas.
\end{acknowledgements}



% ============================================================================
%  BIBLIOGRAFÍA
%  aa.bst debe viajar junto al fuente en el envío (está en class/).
% ============================================================================
\bibliographystyle{aa}
\bibliography{refs}


% ============================================================================
%  APÉNDICES
%
%  Van SIEMPRE después de la bibliografía. Desde julio de 2021 A&A los publica
%  como material camera-ready: el editor no los tipografía, no los corrige y
%  no cambia su layout. Revisar el resultado con cuidado extra.
% ============================================================================
\begin{appendix}
% ============================================================================
%  APÉNDICE A
%
%  ATENCIÓN — Desde julio de 2021, A&A publica los apéndices COMO MATERIAL
%  CAMERA-READY. El editor no los tipografía, no los corrige y no cambia su
%  layout después de recibir la versión aceptada. Lo que compiles aquí es
%  literalmente lo que se publica.
%
%  Numeración: la clase convierte \section en "Appendix A", "Appendix B", ...
%  Las subsecciones quedan A.1, A.2. Referenciar con \appref{app:derivations}.
%
%  FLOATS EN APÉNDICES
%  Los floats deben quedar bajo su propio apéndice: ni antes del título, ni
%  después del título del siguiente. En apéndices cortos, un {figure*} o
%  {table*} genera una página en blanco. Receta recomendada por A&A:
%
%      \onecolumn
%      \section{Título del apéndice}
%      \begin{figure*}[ht!] ... \end{figure*}
%      texto ...
%      \FloatBarrier
%      \twocolumn
%
%  Método alternativo, solo si el apéndice contiene floats y nada de texto:
%  declarar el float primero y meter el \section DENTRO del float.
%
%  Si los floats se siguen escapando de su contexto, usar \onecolumn en todo
%  el bloque de apéndices.
% ============================================================================

\section{Derivations}
\label{app:derivations}

Texto del apéndice.

% ============================================================================
%  APÉNDICE B — Tablas extensas
%
%  Dos vías para una tabla que no cabe en una página:
%
%  1. \aatablelong[n]{csv}{caption}{label} en el CUERPO del artículo.
%     Envuelve en \longtab[n]{...}, que conserva la numeración correlativa y
%     coloca la tabla automáticamente después de los apéndices. Es la vía que
%     A&A recomienda, porque \longtable* solo funciona en onecolumn.
%
%  2. longtable directo dentro de un apéndice en \onecolumn, como abajo.
%
%  Rotadas: \begin{landscape} (paquete lscape) alrededor del longtable, o el
%  entorno sidewaystable* para una única página rotada.
%
%  Las tablas grandes de datos primarios NO van en el artículo: se archivan en
%  el CDS en ASCII con su ReadMe, y en el texto se describe el contenido de
%  cada columna. Ver la skill aa-paper, references/tables.md.
% ============================================================================

\onecolumn
\section{Supplementary tables}
\label{app:tables}

Texto introductorio de las tablas.

% \aatablelong[2]{tables/catalogue.csv}{Full catalogue.}{tab:appcatalogue}

\FloatBarrier
\twocolumn

\end{appendix}

\end{document}
